\section{Análisis del producto: ¿Lovart debe ser herramienta, plataforma o espacio de trabajo?}

Las secciones anteriores trataron el posicionamiento, el flujo de trabajo y las barreras; esta sección analiza Lovart según su forma de producto. Los productos de diseño con IA suelen adoptar tres formas: herramienta, plataforma y espacio de trabajo. Una herramienta resuelve una tarea puntual, como generar un cartel o editar una imagen; una plataforma conecta a creadores, materiales y distribución; un espacio de trabajo sostiene el proceso de producción cotidiano de un rol profesional. Si Lovart quiere convertirse en un Agent vertical, se parece más a un espacio de trabajo que a una herramienta puntual.

\begin{knowledgebox}{Herramienta, plataforma, espacio de trabajo}
\begin{tabularx}{\textwidth}{p{0.20\textwidth}p{0.34\textwidth}X}
\toprule
Forma & Por qué llega el usuario & Origen de la barrera \\
\midrule
Herramienta & Completar rápido una tarea puntual & Experiencia de uso de las funciones y bajo costo. \\
Plataforma & Buscar materiales, plantillas, usuarios o distribución & Efectos de red y oferta de contenido. \\
Espacio de trabajo & Gestionar proyectos, activos, colaboración y entregas & Flujos de trabajo acumulados y costo de migración para la organización. \\
\bottomrule
\end{tabularx}
\end{knowledgebox}

\begin{importantbox}{A Lovart le conviene más ser un espacio de trabajo}
Si Lovart solo es una herramienta, las empresas de modelos la convertirán pronto en una función más; si es una plataforma, tendrá que resolver los problemas de oferta y distribución; si es un espacio de trabajo, puede construir barreras más profundas alrededor del flujo de trabajo continuo del diseñador.
\end{importantbox}

\subsection{Los cuatro ciclos cerrados del espacio de trabajo}

Un espacio de trabajo de diseño necesita cuatro ciclos cerrados: el de proyecto, el de activos, el de retroalimentación y el de entrega. El ciclo de proyecto permite dar seguimiento al brief, las tareas y el avance; el ciclo de activos permite reutilizar los materiales de marca, las imágenes de referencia y los resultados generados; el ciclo de retroalimentación hace que las opiniones del cliente o del equipo entren en la siguiente iteración; el ciclo de entrega hace que los archivos finales sean editables, entregables y de uso comercial.

\begin{knowledgebox}{Tabla de ciclos cerrados del espacio de trabajo de diseño}
\begin{tabularx}{\textwidth}{p{0.20\textwidth}p{0.34\textwidth}X}
\toprule
Ciclo cerrado & Qué resuelve & Consecuencia de no tenerlo \\
\midrule
Ciclo de proyecto & Brief, tareas, avance, versiones & El usuario solo puede hacer generaciones improvisadas. \\
Ciclo de activos & Materiales de marca, propuestas anteriores, imágenes de referencia & Cada vez se empieza de cero y es difícil profesionalizarse. \\
Ciclo de retroalimentación & Comentarios de modificación y registros de revisión & El Agent no entiende las preferencias del equipo. \\
Ciclo de entrega & Formatos de archivo, derechos de autor, tamaños, colaboración & El trabajo no puede entrar en la producción real. \\
\bottomrule
\end{tabularx}
\end{knowledgebox}

\subsection{Resumen de la sección}

El techo del producto Lovart depende de si logra pasar de herramienta a espacio de trabajo. La ventaja defensiva de un Agent vertical no está en una generación aislada, sino en los ciclos cerrados de proyectos y activos a largo plazo.
